\section{Aspekty zdrowotne i środowiskowe}

% SLAJD 11
\begin{frame}{Jakość środowiska naturalnego}
    \begin{itemize}
        \item \textbf{Fundament zdrowia:} Bezpośredni kontakt z nieskażoną przyrodą.
        \item \textbf{Czyste powietrze i cisza:} 
        \begin{itemize}
            \item Redukcja ryzyka chorób cywilizacyjnych (alergie, astma).
            \item Brak permanentnego hałasu miejskiego.
        \end{itemize}
        \item \textbf{Psychika:} Wszechobecna zieleń sprzyja regeneracji psychicznej, redukcji stresu i odnalezieniu wewnętrznej harmonii.
    \end{itemize}
\end{frame}

% SLAJD 12
\begin{frame}{Analiza zanieczyszczeń (PM2.5)}
    \begin{table}
        \centering
        \small
        \caption{Porównanie średniorocznego stężenia pyłu PM2.5 \cite{link1}}
        \begin{tabular}{|l|c|c|}
            \hline
            \textbf{Lokalizacja} & \textbf{Stężenie ($\mu$g/m$^3$)} & \textbf{Norma WHO} \\ \hline
            Duże miasto (Kraków) & 25-35 & 5 \\ \hline
            Wieś (brak przemysłu) & 10-15 & 5 \\ \hline
        \end{tabular}
    \end{table}

    \vspace{0.5cm}
    
    Na wsi stężenie szkodliwego pyłu jest nawet \textbf{2-3 razy niższe}.
    
    \begin{block}{Wzór na stężenie masowe pyłu}
        $$ \text{C}_{\text{PM}_{2.5}} = \frac{m_{\text{PM}_{2.5}}}{V} $$
    \end{block}
\end{frame}

\section{Aspekty społeczne i ekonomiczne}

% SLAJD 13
\begin{frame}{Społeczność i tożsamość}
    \begin{columns}
        \begin{column}{0.5\textwidth}
            \textbf{Więzi społeczne:}
            \begin{itemize}
                \item Brak miejskiej anonimowości.
                \item Relacje oparte na zaufaniu i pomocy sąsiedzkiej (kapitał społeczny).
                \item Wyższe poczucie bezpieczeństwa i oparcia w grupie.
            \end{itemize}
        \end{column}
        
        \begin{column}{0.5\textwidth}
            \textbf{Kultura i tradycja:}
            \begin{itemize}
                \item Żywy kontakt z lokalnym folklorem i obrzędami.
                \item Ochrona przed globalizacją kulturową.
                \item Wzmacnianie tożsamości regionalnej i międzypokoleniowej.
            \end{itemize}
        \end{column}
    \end{columns}
\end{frame}

% SLAJD 14
\begin{frame}{Przestrzeń, wolność i ekonomia}
    \begin{itemize}
        \item \textbf{Ekonomia:} Niższe ceny nieruchomości pozwalają na posiadanie domu z ogrodem (zamiast ciasnego mieszkania).
        \item \textbf{Wolność osobista:} 
        \begin{itemize}
            \item Większa prywatność i brak tłumów.
            \item Możliwość realizacji głośnych hobby czy hodowli zwierząt.
        \end{itemize}
        \item \textbf{Niezależność:} Możliwość częściowej samowystarczalności (własne uprawy) i niezależności finansowej.
    \end{itemize}
\end{frame}

% SLAJD 15
\begin{frame}{Rytm życia i rozwój umiejętności}
    \begin{itemize}
        \item \textbf{Slow Life:} Ucieczka od pośpiechu i presji czasu na rzecz życia zgodnego z rytmem natury.
        \item \textbf{Wartości:} Odchodzenie od konsumpcjonizmu na rzecz autentycznych potrzeb i relacji.
        \item \textbf{Praktyczne umiejętności:}
        \begin{itemize}
            \item Rozwój zaradności (prace w ogrodzie, naprawy).
            \item Poczucie sprawczości, które zanika w miejskim stylu życia.
            \item Szacunek dla zasobów i pracy fizycznej.
        \end{itemize}
    \end{itemize}
\end{frame}